\subsection{Weakened ambient class}

\begin{definition}[Weakened ambient admissibility category]
\label{def:physadm}
Let $\PhysAdmOne$ denote the category whose objects are tuples
\[
\mathfrak P=
(\mathcal A_+,\mathcal H_+,D_+,J,\Gamma,B_\Sigma,\tau_t,\omega)
\]
satisfying the following ambient admissibility axioms:
\begin{enumerate}[leftmargin=2em, topsep=3pt, itemsep=3pt]
    \item[\textbf{(U1)}] \textbf{Real even boundary datum:}
    $(\mathcal A_+,\mathcal H_+,D_+,J,\Gamma,B_\Sigma)$ is an admissible real even one-sided
    boundary datum with reflection-regular collar.
    \item[\textbf{(U4)}] \textbf{Quasilocal positive dynamics:}
    $\tau_t$ is a quasilocal dynamics with reflection positivity and an admissible vacuum sector.
\end{enumerate}
Morphisms are unitary intertwiners preserving the displayed structures.
\end{definition}

No external coverage claim is made in this section. In particular, the ambient class is a theorem
class for internal rigidity, not yet a proof that reality itself is generated from a thinner
primitive layer.

\begin{tcolorbox}[colback=orange!5, colframe=orange!70!black, breakable,
    title={Carrier closure is now split cleanly},
    colbacktitle=orange!75!black]
\Cref{thm:carrier-minimality-k} is now only the algebraic normal form: if the two exterior
signatures
\[
\dim\Lambda^2E_+=1,
\qquad
\Lambda^2E_-\cong E_-^\vee
\]
hold, then one necessarily has
\[
\rank E=5,
\qquad
(\dim E_-,\dim E_+)=(3,2),
\qquad
\dim\Lambda^{\mathrm{even}}E=16.
\]
The internal derivation is completed later by \Cref{cor:bosonic-rank-two,thm:yukawa-forces-32}:
the compact Higgs index fixes $\dim E_+=2$, the primitive indecomposable generator of the actual
branch Yukawa tensor forces $\dim E_-=3$, and the former occupancy law is reduced to a derived
consistency check.
Accordingly $\PhysAdmOne$ no longer carries a packet-size axiom; the rigid minimal carrier is now a
theorem-level closure package rather than a hidden ambient input.
\end{tcolorbox}

\begin{definition}[Internal TFPT reduction functor]
\label{def:universal-reduction}
Define the functor
\[
\Runiv:\PhysAdmOne\to \mathbf{TFPT}^{\mathrm{rig}}_{\mathrm{pre}}
\]
on objects by extracting from $\mathfrak P$ the TFPT pre-class data
\[
\Runiv(\mathfrak P):=
\bigl(
    \Tbound(\mathfrak P),
\Padm(\mathfrak P),
X_f^\circ(\mathfrak P),
[\nabla_F](\mathfrak P),
\Phi(\mathfrak P),
U_\Sigma(\mathfrak P)
\bigr),
\]
where each component is reconstructed by the primitive reconstruction, internal hard-carrier
reduction, internal family reduction, internal transport closure, compact bosonic index, and
seam-transfer closure already
established in the main theorem stack.
\end{definition}

\begin{definition}[Canonical rigid-closure functor]
\label{def:rigid-closure-functor}
Let
\[
\Clrig:\mathbf{TFPT}^{\mathrm{rig}}_{\mathrm{pre}}\to \mathbf{TFPT}^{\mathrm{rig}}
\]
denote the canonical closure functor that adjoins to each pre-class object $\Tpre$ its unique
intrinsic branch $x_{\mathrm{cl}}^\star(\Tpre)$ from \Cref{thm:branch-uniqueness} and then reads
the resulting closed realization in the canonical rigid image via
\Cref{def:tfpt-rigid-category,thm:spectral-completeness-rigid-tfpt}. On morphisms,
$\Clrig$ acts by the same unitary intertwiners. Define the rigid internal-reduction functor by
\[
\Runivrig:=\Clrig\circ \Runiv:\PhysAdmOne\to \mathbf{TFPT}^{\mathrm{rig}}.
\]
\end{definition}

\begin{theorem}[Internal hard-carrier reduction]
\label{prop:universal-carrier-reduction}
For every $\mathfrak P\in\PhysAdmOne$, axiom \textbf{(U1)} together with the boundary-polarization,
boundary-winding, compact Higgs-index, and branch-Yukawa-rigidity theorems reduces functorially to
the unique TFPT hard carrier
\[
E_3\oplus E_2,
\qquad
X=-2P_3+3P_2,
\qquad
Y=-\frac13P_3+\frac12P_2,
\qquad
S^+=\Lambda^{\mathrm{even}}(E_3\oplus E_2).
\]
\end{theorem}

\begin{proof}
Take an ambient object
\[
\mathfrak P=(A_+,H_+,D_+,J,\Gamma,B_\Sigma,\tau_t,\omega)\in\PhysAdmOne.
\]
By axiom \textbf{(U1)}, the tuple $(A_+,H_+,D_+,J,\Gamma,B_\Sigma)$ is an admissible real even
one-sided boundary datum with reflection-regular collar. The boundary-polarization theorem
therefore reconstructs from $\mathfrak P$ the canonical boundary polarization $\iota_C(\mathfrak
P)$ and hence the carrier decomposition on the finite carrier block,
\[
E(\mathfrak P)=E_-(\mathfrak P)\oplus E_+(\mathfrak P).
\]

Next, the boundary-winding theorem determines the canonical unit winding class on the same branch.
The compact Higgs index then fixes the positive carrier rank,
\[
\dim E_+(\mathfrak P)=2.
\]
Once this bosonic rank is known, \Cref{thm:yukawa-forces-32} forces the only admissible negative
rank to be
\[
(\dim E_-,\dim E_+)=(3,2),
\]
and \Cref{thm:carrier-minimality-k,cor:yukawa-discharges-kpp} then recover the former exterior
signatures and packet-size statement as consequences. By the primitive trace condition the
corresponding primitive two-point generator is uniquely
\[
X(\mathfrak P)=-2P_3(\mathfrak P)+3P_2(\mathfrak P),
\]
and therefore
\[
Y(\mathfrak P)=\frac{X(\mathfrak P)}6=-\frac13P_3(\mathfrak P)+\frac12P_2(\mathfrak P).
\]
The exterior packet is then forced to be
\[
S^+(\mathfrak P)=\Lambda^{\mathrm{even}}(E_3(\mathfrak P)\oplus E_2(\mathfrak P)).
\]
Thus the displayed hard-carrier package is determined objectwise.

Now let
\[
u:\mathfrak P\to \mathfrak P'
\]
be a morphism in $\PhysAdmOne$. By definition, $u$ is a unitary intertwiner preserving the
structures listed in Definition~12.1. Therefore it intertwines the boundary operator, the
reflection data, the Calder\'on projector, and the induced boundary polarization. Hence
\[
u\,\iota_C(\mathfrak P)\,u^{-1}=\iota_C(\mathfrak P').
\]
So $u$ carries the carrier block of $\mathfrak P$ unitarily onto the carrier block of $\mathfrak
P'$, preserving the decomposition into the $(-1)$ and $(+1)$ eigenspaces of $\iota_C$.

Because the winding class and the carrier split are uniqueness outputs, the same unitary also
intertwines the derived projectors $P_3,P_2$ and therefore the derived generators $X,Y$. Exterior
functoriality then gives
\[
u\,S^+(\mathfrak P)\,u^{-1}=S^+(\mathfrak P').
\]
Hence the assignment
\[
\mathfrak P\longmapsto (E_3\oplus E_2,X,Y,S^+)
\]
is functorial. This is the asserted internal hard-carrier reduction.
\end{proof}

\begin{theorem}[Internal family reduction]
\label{prop:universal-family-reduction}
For every $\mathfrak P\in\PhysAdmOne$, the boundary-winding theorem and $D_4$-equivariant
monodromy rigidity reduce functorially to
the rigid family data
\[
X_f^\circ\cong \mathbb P^1\setminus\mu_4,
\qquad
[\nabla_F]=[\nabla_F^\star].
\]
\end{theorem}

\begin{proof}
For an ambient object $\mathfrak P\in\PhysAdmOne$, the boundary-winding theorem fixes the
quarter-turn balance of the admissible seam data. The family theorem then identifies the family
section uniquely as the four-punctured sphere
\[
X_f^\circ(\mathfrak P)\cong \mathbb P^1\setminus\mu_4.
\]
The puncture set carries its inherited faithful $D_4$ symmetry.

On that rigid puncture square, the $D_4$-equivariant monodromy-rigidity theorem fixes the
monodromy representation uniquely up to global $SU(3)_F$ conjugation,
\[
\rho_F:\pi_1(X_f^\circ)\to SU(3)_F.
\]
Once the monodromy representation is fixed, the Riemann--Hilbert correspondence for flat Hermitian
bundles yields a unique flat Hermitian family-connection class realizing it. This class is
precisely
\[
[\nabla_F^\star].
\]
Thus the family geometry and holonomy are determined objectwise.

Now let
\[
u:\mathfrak P\to \mathfrak P'
\]
be a morphism. Since $u$ preserves the boundary datum and the induced winding data, it preserves the
quarter-turn balance and hence the puncture-square structure. Therefore it induces an isomorphism
of the two four-punctured sections preserving the $D_4$ action. Under this induced identification
the puncture loops correspond, so the monodromy representations agree up to the same global
$SU(3)_F$ conjugation already allowed by \Cref{thm:d4-monodromy-rigidity}.

Because the flat Hermitian connection class is uniquely determined by that conjugacy class, the
transported connection class coincides with the target one. Therefore
\[
u:(X_f^\circ(\mathfrak P),[\nabla_F(\mathfrak P)])
\longrightarrow
(X_f^\circ(\mathfrak P'),[\nabla_F(\mathfrak P')])
\]
is functorial. This proves the internal family reduction.
\end{proof}

\begin{theorem}[Internal transport reduction]
\label{prop:universal-transport-reduction}
For every $\mathfrak P\in\PhysAdmOne$, the boundary-polarization, carrier-rigidity,
family-holonomy, determinant-line, and positivity theorems reduce functorially to
the TFPT transport package with admissible cusp set
\[
\left\{1,\frac23,\frac13\right\},
\]
exact branch pole $\delta_{\mathrm{ph}}$, and the unique first-generation winding lock.
\end{theorem}

\begin{proof}[Proof sketch]
Once the boundary winding, rigid carrier generator, and rigid family holonomy are fixed,
the admissible cusp spectrum is no longer ambient input; it is read from the singlet-projected
carrier spectrum. The algebraic transport-pole theorem, the exact $C_6$ kernel, and the transport
length sum rule then fix the admissible cusp set, the pole $\delta_{\mathrm{ph}}$, and the
first-generation winding pattern functorially.
\end{proof}

\begin{theorem}[Internal glueing and invariant reconstruction]
\label{prop:universal-glueing}
For every $\mathfrak P\in\PhysAdmOne$, the data extracted in
\Cref{prop:universal-carrier-reduction,prop:universal-family-reduction,prop:universal-transport-reduction}
assemble into a TFPT pre-class object whose invariant tuple is the canonical TFPT tuple
\[
\mathcal I_\star=
\bigl(
    X,Y,[u_\Sigma],X_f^\circ,c_1(L_2),c_1(L_3),[\nabla_F^\star],a_\Sigma,\Lambda_{\mathrm{IR}},x_{\mathrm{cl}}^\star
\bigr).
\]
\end{theorem}

\begin{proof}[Proof sketch]
The carrier, family, transport, determinant-line, gravity, cosmology, and vacuum closures
determine exactly the displayed invariant tuple on the rigid branch. The internal reductions above
therefore glue into one TFPT pre-class object with the same invariant content as the canonical
branch.
\end{proof}

\iffalse
\begin{theorem}[Internal reduction on the weakened ambient class]
\label{thm:universality-tfpt}
The functor
\[
\Runiv:\PhysAdmOne\to \mathbf{TFPT}^{\mathrm{rig}}_{\mathrm{pre}}
\]
is full and faithful. Its rigid closure
\[
\Runivrig=\Clrig\circ \Runiv:\PhysAdmOne\to \mathbf{TFPT}^{\mathrm{rig}}
\]
is essentially surjective. Consequently
\[
\text{within the ambient class }\PhysAdmOne,\qquad
\PhysAdmOne/\!\cong\;\simeq\;\mathbf{TFPT}^{\mathrm{rig}},
\]
and therefore, by \Cref{cor:no-alternatives-tfpt},
\[
\forall \mathfrak P\in\PhysAdmOne:\qquad
\Runivrig(\mathfrak P)\cong \Tstar.
\]
\end{theorem}

\begin{proof}[Proof sketch]
\Cref{prop:universal-carrier-reduction,prop:universal-family-reduction,prop:universal-transport-reduction,prop:universal-glueing}
show that every ambient object of $\PhysAdmOne$ reduces functorially to one rigid TFPT
pre-class object with the canonical invariant tuple. Faithfulness and fullness follow because the
morphisms on both sides are unitary intertwiners of exactly the displayed structures, and the
reconstruction chain forgets no theorem-level invariant. The passage from the pre-class to the
rigid class is the explicit role of \Cref{def:rigid-closure-functor}: \Cref{thm:branch-uniqueness}
supplies the unique closed branch on every pre-class object, and
\Cref{thm:spectral-completeness-rigid-tfpt,cor:no-alternatives-tfpt} collapse the resulting rigid
closure to the unique object $\Tstar$. Hence essential surjectivity is a property of
$\Runivrig=\Clrig\circ\Runiv$, not of $\Runiv$ by itself.
\end{proof}

\fi

\begin{theorem}[Reduction to the canonical rigid orbit]
\label{thm:universality-tfpt}
For every ambient object $\mathfrak P\in\PhysAdmOne$, the rigid internal reduction lands in the
canonical rigid orbit:
\[
\Runivrig(\mathfrak P)\cong \Tstar.
\]
Equivalently, the induced map on isomorphism classes
\[
\pi_0(\PhysAdmOne)\longrightarrow \pi_0(\mathbf{TFPT}^{\mathrm{rig}}),
\qquad
[\mathfrak P]\longmapsto [\Runivrig(\mathfrak P)],
\]
has singleton image
\[
\pi_0(\mathbf{TFPT}^{\mathrm{rig}})=\{[\Tstar]\}.
\]
If $u:\mathfrak P_1\to\mathfrak P_2$ is an ambient morphism, then $\Runivrig(u)$ is the unique rigid
intertwiner between the two reduced objects.
\end{theorem}

\begin{proof}
\Cref{prop:universal-carrier-reduction,prop:universal-family-reduction,prop:universal-transport-reduction,prop:universal-glueing}
produce from every ambient object a TFPT pre-object with the canonical invariant tuple. By
\Cref{thm:branch-uniqueness}, that pre-object carries a unique intrinsic branch, and by
\Cref{thm:spectral-completeness-rigid-tfpt,cor:no-alternatives-tfpt} every rigid closed realization
with this invariant tuple is uniquely isomorphic to $\Tstar$. Therefore every reduced ambient
object lies in the unique rigid orbit of $\Tstar$. The functoriality on morphisms is immediate
because the rigid target has only the unique intertwiner allowed by essential injectivity on the
rigid image.
\end{proof}

\begin{remark}[Primitive honesty]
The manuscript now proves the passage from a minimal operational seed to the one-sided boundary
datum and then the canonical closure from that datum. It does not yet prove that $\Seedop$ itself is
generated from a still thinner universal seed. Accordingly the manuscript establishes internal
rigidity before ultimate primitive minimality.
\end{remark}

\begin{remark}[Internal rigidity versus external adequacy]
The main text proves uniqueness of the closed TFPT object inside the ambient class $\PhysAdmOne$:
\[
\forall \mathfrak P\in\PhysAdmOne,\qquad
\Runivrig(\mathfrak P)\cong \Tstar.
\]
This does not by itself prove that $\PhysAdmOne$ is the physically correct ambient class. External
adequacy is a separate comparison claim on the observable map
\[
\mathcal E:\PhysAdmOne\to\mathcal O_{\mathrm{phys}},
\qquad
\mathfrak P\longmapsto \mathcal E(\mathfrak P).
\]
\end{remark}

The input-side reduction and the output-side orbit are summarized in
\Cref{fig:internal-reduction-orbit}.

\begin{figure}[t]
\centering
\scriptsize
\resizebox{0.94\linewidth}{!}{%
\begin{tikzpicture}[
    >=Latex,
    ambient/.style={draw=black!60, fill=gray!8, rounded corners, text width=2.7cm, minimum height=0.95cm, align=center},
    stage/.style={draw=blue!60!black, fill=blue!6, rounded corners, text width=2.8cm, minimum height=0.95cm, align=center},
    star/.style={draw=green!50!black, fill=green!8, rounded corners, text width=2.7cm, minimum height=0.95cm, align=center},
    orbit/.style={draw=orange!75!black, fill=orange!8, rounded corners, text width=3.2cm, minimum height=0.95cm, align=center},
    band/.style={draw=gray!35, rounded corners, inner sep=10pt}
]
\node[ambient] (phys) at (-6.0,0) {ambient class\\$\PhysAdmOne$};
\node[stage] (pre) at (-2.1,0) {pre-class\\$\mathbf{TFPT}^{\mathrm{rig}}_{\mathrm{pre}}$};
\node[stage] (rig) at (1.8,0) {rigid class\\$\mathbf{TFPT}^{\mathrm{rig}}$};
\node[star] (tstar) at (1.8,-2.0) {$\Tstar$\\unique rigid image};
\node[orbit] (obs) at (6.0,0) {observable hierarchy\\$\GrenTFPT\to\OphysTFPT\to\OschemeTFPT/\SchGrp$};

\node[font=\bfseries, text=tfptblue] at (-2.1,1.8) {input-side internal reduction};
\node[font=\bfseries, text=tfptorange] at (6.0,1.8) {renormalized / physical / scheme readout};

\draw[->, thick] (phys) -- node[above, font=\scriptsize] {$\Runiv$} (pre);
\draw[->, thick] (pre) -- node[above, font=\scriptsize] {$\Clrig$} (rig);
\draw[->, thick] (rig) -- (tstar);
\draw[->, thick, dashed, orange!75!black] (tstar) -- node[below right, font=\scriptsize] {$\Rren,\ \Mphys,\ \Mscheme$} (obs);
\end{tikzpicture}}
\caption{Input-side reduction and observable hierarchy.}
\label{fig:internal-reduction-orbit}
\end{figure}

